\chapter{Fundamentação Teórica}
\label{cap:trabalhos_relacionados}

Em comparação com trabalhos no estado da arte, a nossa arquitetura alinha-se às descobertas de Apostol, Cojocaru e Truică (2024) sobre o uso do ecossistema Spark para grafos de larga escala \cite{apostol2024largescale}. Enquanto esses autores focam-se na escalabilidade de algoritmos de deteção de comunidades (como Louvain e K-Cliques) usando GraphFrames, o presente trabalho garante um passo anterior crucial: a persistência relacional e a modelação da rede de coautoria do Lattes. Além disso, conforme estabelecido no desenho estrutural original do GraphX por Gonzalez \textit{et al.} (2014), a abordagem de manter os atributos dos vértices numa base estruturada (como o PostgreSQL) e delegar a topologia para o Spark minimiza drasticamente o \textit{shuffle} de rede entre os nós de processamento \cite{gonzalez2014graphx}.

A análise de grandes redes complexas, como redes de colaboração acadêmica e social, impulsionam a perspectiva de viabilização de uso de grandes grafos como ferramenta para gerenciar a visualização e relacionamento desses inumeros nós aplicados em ambiente de banco de dados relacionais quando usados em um contexto específico em que bancos de dados nativos como neo4j não atendem a objetivos específicos estipulados.

Ademais, ao verificar o uso de máquina e processamento desses grafos nessa rede complexa, observa-se um desafio para a computação hodierna para lidar com tantos dados em uma só máquina limitada àquele \textit{hardware} específico.

Nesse contexto, é viável a utilização de frameworks que viabilizem a computação distribuida de grafos, como os da distribuição Apache.

No domínio da persistência e processamento de grafos, abordagens iniciais focavam em sistemas nativos de grafos (como Neo4j) operando em instâncias únicas. No entanto, para cenários de \textit{Big Data}, ferramentas baseadas no modelo \textit{MapReduce} e, posteriormente, arquiteturas de processamento em memória como o \textit{Apache Spark} através das bibliotecas \textit{GraphX} e \textit{GraphFrames}, demonstraram superioridade no processamento iterativo de algoritmos topológicos.

Trabalhos de arquitetura híbrida, como os descritos por desenvolvedores em comunidades como o \textit{Stack Overflow}, sugerem o uso do SGBD Relacional como a "fonte da verdade" (\textit{Source of Truth}) para os metadados dos nós, delegando ao \textit{Spark} exclusivamente o processamento intensivo de travessia do grafo. Nossa abordagem alinha-se a essa arquitetura, focando na otimização do fluxo de ida e volta (JDBC -> Spark -> JDBC) dos dados de colaboração científica do sistema Tarrafa.

\textit{Phasellus non faucibus leo. Integer eget varius turpis. Morbi mattis sem sit amet scelerisque consequat. Aenean eget eros ac tortor congue egestas. Nulla pulvinar diam ut odio commodo aliquet.}

Embora existam abordagens que analisem o Currículo Lattes utilizando processamento em lote padrão, a lacuna investigada neste trabalho concentra-se na otimização específica do particionamento e distribuição em um ambiente \textit{Spark} integrado nativamente ao PostgreSQL, visando reduzir a latência na construção e persistência de grandes redes de coautoria.